\documentclass[aspectratio=1610,11pt]{beamer}
% Shared style of the lecture slides (Slides/slidesNN.tex). Each deck starts with
%     \documentclass[aspectratio=1610,11pt]{beamer}
%     % Shared style of the lecture slides (Slides/slidesNN.tex). Each deck starts with
%     \documentclass[aspectratio=1610,11pt]{beamer}
%     % Shared style of the lecture slides (Slides/slidesNN.tex). Each deck starts with
%     \documentclass[aspectratio=1610,11pt]{beamer}
%     \input{slides_preamble}
% aspectratio=1610 (16:10) matches the Samsung Galaxy Tab S9 FE+ (2560 x 1600) used for teaching.
%
% The decks have no overlays (\pause, \only): every slide is one PDF page, so what you write on a
% page while sharing the tablet screen stays on that page. Board pages (\board) are white pages with
% a small header, for calculations and proofs written live.

\usepackage[T1]{fontenc}
\usepackage{lmodern}
\usefonttheme[onlymath]{serif}
\usepackage{amsmath,amssymb}
\usepackage{booktabs}
\usepackage{tikz}
\usetikzlibrary{arrows.meta,calc,positioning,shapes.geometric,decorations.pathreplacing}
\usepackage{pgfplots}
\pgfplotsset{compat=1.18}
\usepgfplotslibrary{fillbetween}

% Numbers of the lecture notes. A deck says which notes it follows with \setnotes{2} and
% loads their labels with \externaldocument[L2-]{../Lecture2/lecture2} (the .aux of the notes
% must exist). Then \begin{theorem}[Spectral theorem][thm:spectral] prints "notes L2, 10.1" at
% the right of the block title, and \secref{sec:spectral} prints "§2.10".
\usepackage{xr-hyper}
\makeatletter
\newcommand{\setnotes}[1]{\def\slides@notes{#1}}
\setnotes{2}
\newcommand{\notesref}[1]{\ref*{L\slides@notes-#1}}
\newcommand{\secref}[1]{\S\slides@notes.\ref*{L\slides@notes-#1}}
\makeatother

% Macros of the lecture notes.
\newcommand{\norm}[1]{\lVert #1 \rVert}
\newcommand{\ip}[2]{\langle #1, #2 \rangle}
\newcommand{\py}[1]{\texttt{\detokenize{#1}}}
\newcommand{\E}{\mathbb{E}}
\newcommand{\indep}{\perp\!\!\!\perp}
\newcommand{\R}{\mathbb{R}}
\newcommand{\bx}{\mathbf{x}}
\newcommand{\by}{\mathbf{y}}
\newcommand{\bw}{\mathbf{w}}
\newcommand{\bv}{\mathbf{v}}
\newcommand{\bu}{\mathbf{u}}
\newcommand{\bq}{\mathbf{q}}
\newcommand{\bb}{\mathbf{b}}
\newcommand{\be}{\mathbf{e}}
\DeclareMathOperator{\Var}{Var}
\DeclareMathOperator{\Cov}{Cov}
\DeclareMathOperator{\Corr}{Corr}
\DeclareMathOperator{\KL}{KL}
\DeclareMathOperator{\im}{im}
\DeclareMathOperator{\rank}{rank}
\DeclareMathOperator{\Null}{null}
\DeclareMathOperator{\col}{col}
\DeclareMathOperator{\row}{row}
\DeclareMathOperator{\Span}{span}
\DeclareMathOperator{\tr}{tr}
\DeclareMathOperator{\diag}{diag}
\DeclareMathOperator{\Div}{div}
\DeclareMathOperator{\curl}{curl}
\DeclareMathOperator{\vol}{vol}
\DeclareMathOperator{\area}{area}
\DeclareMathOperator{\sign}{sign}
\DeclareMathOperator{\softmax}{softmax}
\DeclareMathOperator{\ReLU}{ReLU}
\DeclareMathOperator*{\argmin}{arg\,min}
\DeclareMathOperator*{\argmax}{arg\,max}

% Colors.
\definecolor{accent}{RGB}{31,78,140}      % titles, theorems
\definecolor{defcolor}{RGB}{24,110,92}    % definitions
\definecolor{excolor}{RGB}{150,82,20}     % examples
\definecolor{mlcolor}{RGB}{120,40,120}    % machine learning connections
\definecolor{boardink}{RGB}{140,140,140}  % header of the board pages

% Layout.
\setbeamertemplate{navigation symbols}{}
\setbeamersize{text margin left=0.9em, text margin right=0.9em}
\setbeamercolor{structure}{fg=accent}
\setbeamercolor{normal text}{fg=black, bg=white}
\setbeamercolor{frametitle}{fg=accent, bg=white}
\setbeamerfont{frametitle}{series=\bfseries, size=\large}
\setbeamertemplate{frametitle}{%
    \vskip0.45em\insertframetitle\par\vskip-0.25em
    {\color{accent!25}\rule{\linewidth}{0.5pt}}\vskip-0.2em}
\setbeamertemplate{itemize item}{\color{accent}\raisebox{0.15ex}{\scriptsize$\blacktriangleright$}}
\setbeamertemplate{itemize subitem}{\color{accent}\raisebox{0.2ex}{\tiny$\bullet$}}
\setbeamertemplate{enumerate item}{\color{accent}\insertenumlabel.}
\setbeamerfont{footline}{size=\tiny}
\setbeamertemplate{footline}{%
    \leavevmode\hbox to \paperwidth{\hskip0.9em{\color{gray}\insertshorttitle}\hfill
    {\color{gray}\insertframenumber}\hskip0.9em}\vskip0.45em}
\setbeamertemplate{blocks}[rounded][shadow=false]
\setbeamercolor{block title}{fg=white, bg=accent}
\setbeamercolor{block body}{fg=black, bg=accent!6}
\setbeamerfont{block title}{size=\normalsize, series=\bfseries}
\setlength{\abovedisplayskip}{4pt}
\setlength{\belowdisplayskip}{4pt}

% Blocks: \begin{theorem}[Spectral theorem][thm:spectral] ... \end{theorem}; both arguments are optional.
\makeatletter
\newcommand{\slides@block}[5]{% color, word, name, label in the notes, body tint
    \setbeamercolor{block title}{fg=white, bg=#1}%
    \setbeamercolor{block body}{fg=black, bg=#1!#5}%
    \begin{block}{#2\if\relax\detokenize{#3}\relax\else\ {\normalfont(#3)}\fi
        \if\relax\detokenize{#4}\relax\else\hfill{\normalfont\scriptsize notes~L\slides@notes, \notesref{#4}}\fi}}
\DeclareDocumentEnvironment{theorem}{O{} O{}}{\slides@block{accent}{Theorem}{#1}{#2}{6}}{\end{block}}
\DeclareDocumentEnvironment{proposition}{O{} O{}}{\slides@block{accent}{Proposition}{#1}{#2}{6}}{\end{block}}
\DeclareDocumentEnvironment{lemma}{O{} O{}}{\slides@block{accent}{Lemma}{#1}{#2}{6}}{\end{block}}
\DeclareDocumentEnvironment{corollary}{O{} O{}}{\slides@block{accent}{Corollary}{#1}{#2}{6}}{\end{block}}
\DeclareDocumentEnvironment{definition}{O{} O{}}{\slides@block{defcolor}{Definition}{#1}{#2}{6}}{\end{block}}
\DeclareDocumentEnvironment{example}{O{} O{}}{\slides@block{excolor}{Example}{#1}{#2}{6}}{\end{block}}
\DeclareDocumentEnvironment{remark}{O{} O{}}{\slides@block{gray!70!black}{Remark}{#1}{#2}{5}}{\end{block}}
\DeclareDocumentEnvironment{mlbox}{O{In machine learning}}{%
    \setbeamercolor{block title}{fg=white, bg=mlcolor}%
    \setbeamercolor{block body}{fg=black, bg=mlcolor!6}%
    \begin{block}{#1}}{\end{block}}
\makeatother

% Board pages.
%   \board{Title}{Task, shown in small print at the top}            one white page
%   \board[3]{Title}{Task}                                          three pages (the last two blank)
\newcounter{boardpage}
\newcommand{\boardhead}[1]{%
    \vskip0.3em
    \noindent{\color{boardink}\scriptsize
        \tikz[baseline=(b.base)]\node[draw=boardink, rounded corners=2pt, inner sep=1.5pt,
            font=\tiny\bfseries\sffamily] (b) {BOARD};\ \ #1}\par\vskip0.3em}
\newcommand{\board}[3][1]{%
    \begin{frame}[t]
        \boardhead{#2}
        \if\relax\detokenize{#3}\relax\else
            {\small #3\par}
        \fi
    \end{frame}%
    \setcounter{boardpage}{1}%
    \loop\ifnum\value{boardpage}<#1\relax
        \stepcounter{boardpage}%
        \begin{frame}[t]\boardhead{#2 (continued)}\end{frame}%
    \repeat}

% Practice pages: problems of the practice file, as white pages for solving them live.
% A deck loads the practice file's labels with \externaldocument[P1-]{../Lecture1/practice1} and
% \setpractice{1}. Then
%   \practice{pr:eig2}{\mandatory}{Problem text}    one page, headed "PRACTICE  Problem 10.1 (mandatory)"
%   \practice[2]{pr:eig-props}{}{Problem text}      two pages (the second blank)
%   \practice{pr:eig-gre}{$\star$}{Problem text}    the marks of the practice file
% The answers stay in LectureN/practiceN_answers.tex.
\newcommand{\mandatory}{\textcolor{red!70!black}{\textbf{mandatory}}}
\makeatletter
\newcommand{\setpractice}[1]{\def\slides@practice{#1}}
\newcommand{\practiceref}[1]{\ref*{P\slides@practice-#1}}
\newcommand{\practicehead}[2]{% label in the practice file, text after the problem number
    \vskip0.3em
    \noindent{\color{boardink}\scriptsize
        \tikz[baseline=(b.base)]\node[draw=defcolor, text=defcolor, rounded corners=2pt, inner sep=1.5pt,
            font=\tiny\bfseries\sffamily] (b) {PRACTICE};\ \ Practice~\slides@practice, Problem~\practiceref{#1}#2}%
    \par\vskip0.3em}
\newcommand{\practice}[4][1]{%
    \begin{frame}[t]
        \practicehead{#2}{\if\relax\detokenize{#3}\relax\else\ (#3)\fi}
        {\small #4\par}
    \end{frame}%
    \setcounter{boardpage}{1}%
    \loop\ifnum\value{boardpage}<#1\relax
        \stepcounter{boardpage}%
        \begin{frame}[t]\practicehead{#2}{\ (continued)}\end{frame}%
    \repeat}
\makeatother

% Divider page: \topic{\secref{sec:spectral}}{The spectral theorem} or \topic{\S2.1--2.2}{Norms}.
\newcommand{\topic}[2]{%
    \begin{frame}[c]
        \centering
        {\color{gray}\large #1}\par\vskip0.4em
        {\color{accent}\LARGE\bfseries #2}\par
    \end{frame}}

% Title page.
\setbeamercolor{title}{fg=accent}
\setbeamerfont{title}{series=\bfseries, size=\Large}
\setbeamercolor{subtitle}{fg=gray!60!black}
\setbeamerfont{subtitle}{size=\normalsize}
\setbeamertemplate{title page}{%
    \vfill
    {\usebeamercolor[fg]{title}\usebeamerfont{title}\inserttitle\par}
    \vskip0.6em
    {\usebeamercolor[fg]{subtitle}\usebeamerfont{subtitle}\insertsubtitle\par}
    \vskip1.2em
    {\small\color{gray}\insertauthor\par}
    \vfill}
\author{Mathematics for Machine Learning}
\date{}

% Two columns with a picture: \sidebyside{0.5}{left text}{right picture}.
\newcommand{\sidebyside}[3]{%
    \begin{columns}[T]
        \begin{column}{#1\linewidth}#2\end{column}
        \begin{column}{\dimexpr\linewidth-#1\linewidth-1em\relax}#3\end{column}
    \end{columns}}

% "Check" slide after a board computation.
\newcommand{\checkslide}[2]{%
    \begin{frame}{Check: #1}
        #2
    \end{frame}}

% Pictures.
\newcommand{\notesfig}[2][1]{\resizebox{#1\linewidth}{!}{\input{figures/notes/#2}}}
\newcommand{\fig}[2][1]{\includegraphics[width=#1\linewidth]{figures/#2}}

% aspectratio=1610 (16:10) matches the Samsung Galaxy Tab S9 FE+ (2560 x 1600) used for teaching.
%
% The decks have no overlays (\pause, \only): every slide is one PDF page, so what you write on a
% page while sharing the tablet screen stays on that page. Board pages (\board) are white pages with
% a small header, for calculations and proofs written live.

\usepackage[T1]{fontenc}
\usepackage{lmodern}
\usefonttheme[onlymath]{serif}
\usepackage{amsmath,amssymb}
\usepackage{booktabs}
\usepackage{tikz}
\usetikzlibrary{arrows.meta,calc,positioning,shapes.geometric,decorations.pathreplacing}
\usepackage{pgfplots}
\pgfplotsset{compat=1.18}
\usepgfplotslibrary{fillbetween}

% Numbers of the lecture notes. A deck says which notes it follows with \setnotes{2} and
% loads their labels with \externaldocument[L2-]{../Lecture2/lecture2} (the .aux of the notes
% must exist). Then \begin{theorem}[Spectral theorem][thm:spectral] prints "notes L2, 10.1" at
% the right of the block title, and \secref{sec:spectral} prints "§2.10".
\usepackage{xr-hyper}
\makeatletter
\newcommand{\setnotes}[1]{\def\slides@notes{#1}}
\setnotes{2}
\newcommand{\notesref}[1]{\ref*{L\slides@notes-#1}}
\newcommand{\secref}[1]{\S\slides@notes.\ref*{L\slides@notes-#1}}
\makeatother

% Macros of the lecture notes.
\newcommand{\norm}[1]{\lVert #1 \rVert}
\newcommand{\ip}[2]{\langle #1, #2 \rangle}
\newcommand{\py}[1]{\texttt{\detokenize{#1}}}
\newcommand{\E}{\mathbb{E}}
\newcommand{\indep}{\perp\!\!\!\perp}
\newcommand{\R}{\mathbb{R}}
\newcommand{\bx}{\mathbf{x}}
\newcommand{\by}{\mathbf{y}}
\newcommand{\bw}{\mathbf{w}}
\newcommand{\bv}{\mathbf{v}}
\newcommand{\bu}{\mathbf{u}}
\newcommand{\bq}{\mathbf{q}}
\newcommand{\bb}{\mathbf{b}}
\newcommand{\be}{\mathbf{e}}
\DeclareMathOperator{\Var}{Var}
\DeclareMathOperator{\Cov}{Cov}
\DeclareMathOperator{\Corr}{Corr}
\DeclareMathOperator{\KL}{KL}
\DeclareMathOperator{\im}{im}
\DeclareMathOperator{\rank}{rank}
\DeclareMathOperator{\Null}{null}
\DeclareMathOperator{\col}{col}
\DeclareMathOperator{\row}{row}
\DeclareMathOperator{\Span}{span}
\DeclareMathOperator{\tr}{tr}
\DeclareMathOperator{\diag}{diag}
\DeclareMathOperator{\Div}{div}
\DeclareMathOperator{\curl}{curl}
\DeclareMathOperator{\vol}{vol}
\DeclareMathOperator{\area}{area}
\DeclareMathOperator{\sign}{sign}
\DeclareMathOperator{\softmax}{softmax}
\DeclareMathOperator{\ReLU}{ReLU}
\DeclareMathOperator*{\argmin}{arg\,min}
\DeclareMathOperator*{\argmax}{arg\,max}

% Colors.
\definecolor{accent}{RGB}{31,78,140}      % titles, theorems
\definecolor{defcolor}{RGB}{24,110,92}    % definitions
\definecolor{excolor}{RGB}{150,82,20}     % examples
\definecolor{mlcolor}{RGB}{120,40,120}    % machine learning connections
\definecolor{boardink}{RGB}{140,140,140}  % header of the board pages

% Layout.
\setbeamertemplate{navigation symbols}{}
\setbeamersize{text margin left=0.9em, text margin right=0.9em}
\setbeamercolor{structure}{fg=accent}
\setbeamercolor{normal text}{fg=black, bg=white}
\setbeamercolor{frametitle}{fg=accent, bg=white}
\setbeamerfont{frametitle}{series=\bfseries, size=\large}
\setbeamertemplate{frametitle}{%
    \vskip0.45em\insertframetitle\par\vskip-0.25em
    {\color{accent!25}\rule{\linewidth}{0.5pt}}\vskip-0.2em}
\setbeamertemplate{itemize item}{\color{accent}\raisebox{0.15ex}{\scriptsize$\blacktriangleright$}}
\setbeamertemplate{itemize subitem}{\color{accent}\raisebox{0.2ex}{\tiny$\bullet$}}
\setbeamertemplate{enumerate item}{\color{accent}\insertenumlabel.}
\setbeamerfont{footline}{size=\tiny}
\setbeamertemplate{footline}{%
    \leavevmode\hbox to \paperwidth{\hskip0.9em{\color{gray}\insertshorttitle}\hfill
    {\color{gray}\insertframenumber}\hskip0.9em}\vskip0.45em}
\setbeamertemplate{blocks}[rounded][shadow=false]
\setbeamercolor{block title}{fg=white, bg=accent}
\setbeamercolor{block body}{fg=black, bg=accent!6}
\setbeamerfont{block title}{size=\normalsize, series=\bfseries}
\setlength{\abovedisplayskip}{4pt}
\setlength{\belowdisplayskip}{4pt}

% Blocks: \begin{theorem}[Spectral theorem][thm:spectral] ... \end{theorem}; both arguments are optional.
\makeatletter
\newcommand{\slides@block}[5]{% color, word, name, label in the notes, body tint
    \setbeamercolor{block title}{fg=white, bg=#1}%
    \setbeamercolor{block body}{fg=black, bg=#1!#5}%
    \begin{block}{#2\if\relax\detokenize{#3}\relax\else\ {\normalfont(#3)}\fi
        \if\relax\detokenize{#4}\relax\else\hfill{\normalfont\scriptsize notes~L\slides@notes, \notesref{#4}}\fi}}
\DeclareDocumentEnvironment{theorem}{O{} O{}}{\slides@block{accent}{Theorem}{#1}{#2}{6}}{\end{block}}
\DeclareDocumentEnvironment{proposition}{O{} O{}}{\slides@block{accent}{Proposition}{#1}{#2}{6}}{\end{block}}
\DeclareDocumentEnvironment{lemma}{O{} O{}}{\slides@block{accent}{Lemma}{#1}{#2}{6}}{\end{block}}
\DeclareDocumentEnvironment{corollary}{O{} O{}}{\slides@block{accent}{Corollary}{#1}{#2}{6}}{\end{block}}
\DeclareDocumentEnvironment{definition}{O{} O{}}{\slides@block{defcolor}{Definition}{#1}{#2}{6}}{\end{block}}
\DeclareDocumentEnvironment{example}{O{} O{}}{\slides@block{excolor}{Example}{#1}{#2}{6}}{\end{block}}
\DeclareDocumentEnvironment{remark}{O{} O{}}{\slides@block{gray!70!black}{Remark}{#1}{#2}{5}}{\end{block}}
\DeclareDocumentEnvironment{mlbox}{O{In machine learning}}{%
    \setbeamercolor{block title}{fg=white, bg=mlcolor}%
    \setbeamercolor{block body}{fg=black, bg=mlcolor!6}%
    \begin{block}{#1}}{\end{block}}
\makeatother

% Board pages.
%   \board{Title}{Task, shown in small print at the top}            one white page
%   \board[3]{Title}{Task}                                          three pages (the last two blank)
\newcounter{boardpage}
\newcommand{\boardhead}[1]{%
    \vskip0.3em
    \noindent{\color{boardink}\scriptsize
        \tikz[baseline=(b.base)]\node[draw=boardink, rounded corners=2pt, inner sep=1.5pt,
            font=\tiny\bfseries\sffamily] (b) {BOARD};\ \ #1}\par\vskip0.3em}
\newcommand{\board}[3][1]{%
    \begin{frame}[t]
        \boardhead{#2}
        \if\relax\detokenize{#3}\relax\else
            {\small #3\par}
        \fi
    \end{frame}%
    \setcounter{boardpage}{1}%
    \loop\ifnum\value{boardpage}<#1\relax
        \stepcounter{boardpage}%
        \begin{frame}[t]\boardhead{#2 (continued)}\end{frame}%
    \repeat}

% Practice pages: problems of the practice file, as white pages for solving them live.
% A deck loads the practice file's labels with \externaldocument[P1-]{../Lecture1/practice1} and
% \setpractice{1}. Then
%   \practice{pr:eig2}{\mandatory}{Problem text}    one page, headed "PRACTICE  Problem 10.1 (mandatory)"
%   \practice[2]{pr:eig-props}{}{Problem text}      two pages (the second blank)
%   \practice{pr:eig-gre}{$\star$}{Problem text}    the marks of the practice file
% The answers stay in LectureN/practiceN_answers.tex.
\newcommand{\mandatory}{\textcolor{red!70!black}{\textbf{mandatory}}}
\makeatletter
\newcommand{\setpractice}[1]{\def\slides@practice{#1}}
\newcommand{\practiceref}[1]{\ref*{P\slides@practice-#1}}
\newcommand{\practicehead}[2]{% label in the practice file, text after the problem number
    \vskip0.3em
    \noindent{\color{boardink}\scriptsize
        \tikz[baseline=(b.base)]\node[draw=defcolor, text=defcolor, rounded corners=2pt, inner sep=1.5pt,
            font=\tiny\bfseries\sffamily] (b) {PRACTICE};\ \ Practice~\slides@practice, Problem~\practiceref{#1}#2}%
    \par\vskip0.3em}
\newcommand{\practice}[4][1]{%
    \begin{frame}[t]
        \practicehead{#2}{\if\relax\detokenize{#3}\relax\else\ (#3)\fi}
        {\small #4\par}
    \end{frame}%
    \setcounter{boardpage}{1}%
    \loop\ifnum\value{boardpage}<#1\relax
        \stepcounter{boardpage}%
        \begin{frame}[t]\practicehead{#2}{\ (continued)}\end{frame}%
    \repeat}
\makeatother

% Divider page: \topic{\secref{sec:spectral}}{The spectral theorem} or \topic{\S2.1--2.2}{Norms}.
\newcommand{\topic}[2]{%
    \begin{frame}[c]
        \centering
        {\color{gray}\large #1}\par\vskip0.4em
        {\color{accent}\LARGE\bfseries #2}\par
    \end{frame}}

% Title page.
\setbeamercolor{title}{fg=accent}
\setbeamerfont{title}{series=\bfseries, size=\Large}
\setbeamercolor{subtitle}{fg=gray!60!black}
\setbeamerfont{subtitle}{size=\normalsize}
\setbeamertemplate{title page}{%
    \vfill
    {\usebeamercolor[fg]{title}\usebeamerfont{title}\inserttitle\par}
    \vskip0.6em
    {\usebeamercolor[fg]{subtitle}\usebeamerfont{subtitle}\insertsubtitle\par}
    \vskip1.2em
    {\small\color{gray}\insertauthor\par}
    \vfill}
\author{Mathematics for Machine Learning}
\date{}

% Two columns with a picture: \sidebyside{0.5}{left text}{right picture}.
\newcommand{\sidebyside}[3]{%
    \begin{columns}[T]
        \begin{column}{#1\linewidth}#2\end{column}
        \begin{column}{\dimexpr\linewidth-#1\linewidth-1em\relax}#3\end{column}
    \end{columns}}

% "Check" slide after a board computation.
\newcommand{\checkslide}[2]{%
    \begin{frame}{Check: #1}
        #2
    \end{frame}}

% Pictures.
\newcommand{\notesfig}[2][1]{\resizebox{#1\linewidth}{!}{\input{figures/notes/#2}}}
\newcommand{\fig}[2][1]{\includegraphics[width=#1\linewidth]{figures/#2}}

% aspectratio=1610 (16:10) matches the Samsung Galaxy Tab S9 FE+ (2560 x 1600) used for teaching.
%
% The decks have no overlays (\pause, \only): every slide is one PDF page, so what you write on a
% page while sharing the tablet screen stays on that page. Board pages (\board) are white pages with
% a small header, for calculations and proofs written live.

\usepackage[T1]{fontenc}
\usepackage{lmodern}
\usefonttheme[onlymath]{serif}
\usepackage{amsmath,amssymb}
\usepackage{booktabs}
\usepackage{tikz}
\usetikzlibrary{arrows.meta,calc,positioning,shapes.geometric,decorations.pathreplacing}
\usepackage{pgfplots}
\pgfplotsset{compat=1.18}
\usepgfplotslibrary{fillbetween}

% Numbers of the lecture notes. A deck says which notes it follows with \setnotes{2} and
% loads their labels with \externaldocument[L2-]{../Lecture2/lecture2} (the .aux of the notes
% must exist). Then \begin{theorem}[Spectral theorem][thm:spectral] prints "notes L2, 10.1" at
% the right of the block title, and \secref{sec:spectral} prints "§2.10".
\usepackage{xr-hyper}
\makeatletter
\newcommand{\setnotes}[1]{\def\slides@notes{#1}}
\setnotes{2}
\newcommand{\notesref}[1]{\ref*{L\slides@notes-#1}}
\newcommand{\secref}[1]{\S\slides@notes.\ref*{L\slides@notes-#1}}
\makeatother

% Macros of the lecture notes.
\newcommand{\norm}[1]{\lVert #1 \rVert}
\newcommand{\ip}[2]{\langle #1, #2 \rangle}
\newcommand{\py}[1]{\texttt{\detokenize{#1}}}
\newcommand{\E}{\mathbb{E}}
\newcommand{\indep}{\perp\!\!\!\perp}
\newcommand{\R}{\mathbb{R}}
\newcommand{\bx}{\mathbf{x}}
\newcommand{\by}{\mathbf{y}}
\newcommand{\bw}{\mathbf{w}}
\newcommand{\bv}{\mathbf{v}}
\newcommand{\bu}{\mathbf{u}}
\newcommand{\bq}{\mathbf{q}}
\newcommand{\bb}{\mathbf{b}}
\newcommand{\be}{\mathbf{e}}
\DeclareMathOperator{\Var}{Var}
\DeclareMathOperator{\Cov}{Cov}
\DeclareMathOperator{\Corr}{Corr}
\DeclareMathOperator{\KL}{KL}
\DeclareMathOperator{\im}{im}
\DeclareMathOperator{\rank}{rank}
\DeclareMathOperator{\Null}{null}
\DeclareMathOperator{\col}{col}
\DeclareMathOperator{\row}{row}
\DeclareMathOperator{\Span}{span}
\DeclareMathOperator{\tr}{tr}
\DeclareMathOperator{\diag}{diag}
\DeclareMathOperator{\Div}{div}
\DeclareMathOperator{\curl}{curl}
\DeclareMathOperator{\vol}{vol}
\DeclareMathOperator{\area}{area}
\DeclareMathOperator{\sign}{sign}
\DeclareMathOperator{\softmax}{softmax}
\DeclareMathOperator{\ReLU}{ReLU}
\DeclareMathOperator*{\argmin}{arg\,min}
\DeclareMathOperator*{\argmax}{arg\,max}

% Colors.
\definecolor{accent}{RGB}{31,78,140}      % titles, theorems
\definecolor{defcolor}{RGB}{24,110,92}    % definitions
\definecolor{excolor}{RGB}{150,82,20}     % examples
\definecolor{mlcolor}{RGB}{120,40,120}    % machine learning connections
\definecolor{boardink}{RGB}{140,140,140}  % header of the board pages

% Layout.
\setbeamertemplate{navigation symbols}{}
\setbeamersize{text margin left=0.9em, text margin right=0.9em}
\setbeamercolor{structure}{fg=accent}
\setbeamercolor{normal text}{fg=black, bg=white}
\setbeamercolor{frametitle}{fg=accent, bg=white}
\setbeamerfont{frametitle}{series=\bfseries, size=\large}
\setbeamertemplate{frametitle}{%
    \vskip0.45em\insertframetitle\par\vskip-0.25em
    {\color{accent!25}\rule{\linewidth}{0.5pt}}\vskip-0.2em}
\setbeamertemplate{itemize item}{\color{accent}\raisebox{0.15ex}{\scriptsize$\blacktriangleright$}}
\setbeamertemplate{itemize subitem}{\color{accent}\raisebox{0.2ex}{\tiny$\bullet$}}
\setbeamertemplate{enumerate item}{\color{accent}\insertenumlabel.}
\setbeamerfont{footline}{size=\tiny}
\setbeamertemplate{footline}{%
    \leavevmode\hbox to \paperwidth{\hskip0.9em{\color{gray}\insertshorttitle}\hfill
    {\color{gray}\insertframenumber}\hskip0.9em}\vskip0.45em}
\setbeamertemplate{blocks}[rounded][shadow=false]
\setbeamercolor{block title}{fg=white, bg=accent}
\setbeamercolor{block body}{fg=black, bg=accent!6}
\setbeamerfont{block title}{size=\normalsize, series=\bfseries}
\setlength{\abovedisplayskip}{4pt}
\setlength{\belowdisplayskip}{4pt}

% Blocks: \begin{theorem}[Spectral theorem][thm:spectral] ... \end{theorem}; both arguments are optional.
\makeatletter
\newcommand{\slides@block}[5]{% color, word, name, label in the notes, body tint
    \setbeamercolor{block title}{fg=white, bg=#1}%
    \setbeamercolor{block body}{fg=black, bg=#1!#5}%
    \begin{block}{#2\if\relax\detokenize{#3}\relax\else\ {\normalfont(#3)}\fi
        \if\relax\detokenize{#4}\relax\else\hfill{\normalfont\scriptsize notes~L\slides@notes, \notesref{#4}}\fi}}
\DeclareDocumentEnvironment{theorem}{O{} O{}}{\slides@block{accent}{Theorem}{#1}{#2}{6}}{\end{block}}
\DeclareDocumentEnvironment{proposition}{O{} O{}}{\slides@block{accent}{Proposition}{#1}{#2}{6}}{\end{block}}
\DeclareDocumentEnvironment{lemma}{O{} O{}}{\slides@block{accent}{Lemma}{#1}{#2}{6}}{\end{block}}
\DeclareDocumentEnvironment{corollary}{O{} O{}}{\slides@block{accent}{Corollary}{#1}{#2}{6}}{\end{block}}
\DeclareDocumentEnvironment{definition}{O{} O{}}{\slides@block{defcolor}{Definition}{#1}{#2}{6}}{\end{block}}
\DeclareDocumentEnvironment{example}{O{} O{}}{\slides@block{excolor}{Example}{#1}{#2}{6}}{\end{block}}
\DeclareDocumentEnvironment{remark}{O{} O{}}{\slides@block{gray!70!black}{Remark}{#1}{#2}{5}}{\end{block}}
\DeclareDocumentEnvironment{mlbox}{O{In machine learning}}{%
    \setbeamercolor{block title}{fg=white, bg=mlcolor}%
    \setbeamercolor{block body}{fg=black, bg=mlcolor!6}%
    \begin{block}{#1}}{\end{block}}
\makeatother

% Board pages.
%   \board{Title}{Task, shown in small print at the top}            one white page
%   \board[3]{Title}{Task}                                          three pages (the last two blank)
\newcounter{boardpage}
\newcommand{\boardhead}[1]{%
    \vskip0.3em
    \noindent{\color{boardink}\scriptsize
        \tikz[baseline=(b.base)]\node[draw=boardink, rounded corners=2pt, inner sep=1.5pt,
            font=\tiny\bfseries\sffamily] (b) {BOARD};\ \ #1}\par\vskip0.3em}
\newcommand{\board}[3][1]{%
    \begin{frame}[t]
        \boardhead{#2}
        \if\relax\detokenize{#3}\relax\else
            {\small #3\par}
        \fi
    \end{frame}%
    \setcounter{boardpage}{1}%
    \loop\ifnum\value{boardpage}<#1\relax
        \stepcounter{boardpage}%
        \begin{frame}[t]\boardhead{#2 (continued)}\end{frame}%
    \repeat}

% Practice pages: problems of the practice file, as white pages for solving them live.
% A deck loads the practice file's labels with \externaldocument[P1-]{../Lecture1/practice1} and
% \setpractice{1}. Then
%   \practice{pr:eig2}{\mandatory}{Problem text}    one page, headed "PRACTICE  Problem 10.1 (mandatory)"
%   \practice[2]{pr:eig-props}{}{Problem text}      two pages (the second blank)
%   \practice{pr:eig-gre}{$\star$}{Problem text}    the marks of the practice file
% The answers stay in LectureN/practiceN_answers.tex.
\newcommand{\mandatory}{\textcolor{red!70!black}{\textbf{mandatory}}}
\makeatletter
\newcommand{\setpractice}[1]{\def\slides@practice{#1}}
\newcommand{\practiceref}[1]{\ref*{P\slides@practice-#1}}
\newcommand{\practicehead}[2]{% label in the practice file, text after the problem number
    \vskip0.3em
    \noindent{\color{boardink}\scriptsize
        \tikz[baseline=(b.base)]\node[draw=defcolor, text=defcolor, rounded corners=2pt, inner sep=1.5pt,
            font=\tiny\bfseries\sffamily] (b) {PRACTICE};\ \ Practice~\slides@practice, Problem~\practiceref{#1}#2}%
    \par\vskip0.3em}
\newcommand{\practice}[4][1]{%
    \begin{frame}[t]
        \practicehead{#2}{\if\relax\detokenize{#3}\relax\else\ (#3)\fi}
        {\small #4\par}
    \end{frame}%
    \setcounter{boardpage}{1}%
    \loop\ifnum\value{boardpage}<#1\relax
        \stepcounter{boardpage}%
        \begin{frame}[t]\practicehead{#2}{\ (continued)}\end{frame}%
    \repeat}
\makeatother

% Divider page: \topic{\secref{sec:spectral}}{The spectral theorem} or \topic{\S2.1--2.2}{Norms}.
\newcommand{\topic}[2]{%
    \begin{frame}[c]
        \centering
        {\color{gray}\large #1}\par\vskip0.4em
        {\color{accent}\LARGE\bfseries #2}\par
    \end{frame}}

% Title page.
\setbeamercolor{title}{fg=accent}
\setbeamerfont{title}{series=\bfseries, size=\Large}
\setbeamercolor{subtitle}{fg=gray!60!black}
\setbeamerfont{subtitle}{size=\normalsize}
\setbeamertemplate{title page}{%
    \vfill
    {\usebeamercolor[fg]{title}\usebeamerfont{title}\inserttitle\par}
    \vskip0.6em
    {\usebeamercolor[fg]{subtitle}\usebeamerfont{subtitle}\insertsubtitle\par}
    \vskip1.2em
    {\small\color{gray}\insertauthor\par}
    \vfill}
\author{Mathematics for Machine Learning}
\date{}

% Two columns with a picture: \sidebyside{0.5}{left text}{right picture}.
\newcommand{\sidebyside}[3]{%
    \begin{columns}[T]
        \begin{column}{#1\linewidth}#2\end{column}
        \begin{column}{\dimexpr\linewidth-#1\linewidth-1em\relax}#3\end{column}
    \end{columns}}

% "Check" slide after a board computation.
\newcommand{\checkslide}[2]{%
    \begin{frame}{Check: #1}
        #2
    \end{frame}}

% Pictures.
\newcommand{\notesfig}[2][1]{\resizebox{#1\linewidth}{!}{\input{figures/notes/#2}}}
\newcommand{\fig}[2][1]{\includegraphics[width=#1\linewidth]{figures/#2}}

\externaldocument[L3-]{../Lecture3/lecture3}
\setnotes{3}

\title[Lecture 8 · Calculus II and Optimization]{Lecture 8. Calculus II and Optimization}
\subtitle{Notes: Lecture 3, \S3.9--3.12}

\begin{document}

\begin{frame}[plain]
    \titlepage
\end{frame}

\begin{frame}{Today}
    \begin{columns}[T]
        \begin{column}{0.52\linewidth}
            \begin{enumerate}\small
                \item Critical points and the second-order test
                \item Gradient descent: smooth functions, line search, quadratics and the condition number
                \item Momentum and Newton's method
                \item Convexity, strong convexity, convergence rates
                \item Subgradients, the proximal gradient method, the lasso
                \item Lagrange multipliers, KKT, duality, convex conjugates
            \end{enumerate}
        \end{column}
        \begin{column}{0.44\linewidth}
            \begin{mlbox}[Used in]
                \small
                \begin{itemize}
                    \item GD and SGD for linear and logistic regression (ML)
                    \item L1 and L2 regularization as constrained problems (ML)
                    \item momentum, RMSProp, Adam (DL)
                    \item trust regions in TRPO (RL)
                \end{itemize}
            \end{mlbox}
        \end{column}
    \end{columns}
\end{frame}

% ------------------------------------------------------------------ Critical points
\topic{\secref{sec:optimization}}{Unconstrained optimization}

\begin{frame}{Minima and critical points}
    \begin{definition}[][def:minimum]
        \small $\bx^*$ is a \textbf{global minimum} of $f$ if $f(\bx^*) \leq f(\bx)$ for all $\bx$, a \textbf{local minimum} if this holds for all $\bx$ near $\bx^*$, and \textbf{strict} if $<$ for $\bx \neq \bx^*$. \ A \textbf{critical point} has $\nabla f(\bx^*) = \mathbf{0}$.
    \end{definition}
    \begin{theorem}[First-order condition][thm:first-order]
        \small If $f$ is differentiable and $\bx^*$ is a local minimum, then $\nabla f(\bx^*) = \mathbf{0}$.
    \end{theorem}
    \small
    Proof: $g(t) = f(\bx^* + t\bv)$ has a local minimum at $t = 0$, so $0 = g'(0) = \nabla f(\bx^*)^\top\bv$ for every $\bv$.
    \\[0.5em]
    Critical points can be minima, maxima or saddle points. In high dimensions, most critical points of a random-looking function are saddle points.
\end{frame}

\begin{frame}{The second-order test}
    \begin{theorem}[Second-order test][thm:second-order]
        \small Let $\nabla f(\bx^*) = \mathbf{0}$ and let $\lambda_1 \geq \cdots \geq \lambda_n$ be the eigenvalues of $\nabla^2 f(\bx^*)$.
        \begin{enumerate}
            \item all $\lambda_i > 0$: strict local minimum; \quad all $\lambda_i < 0$: strict local maximum;
            \item $\lambda_1 > 0 > \lambda_n$: saddle point ($f$ increases along eigenvectors with $\lambda > 0$, decreases along those with $\lambda < 0$);
            \item conversely, at a local minimum all $\lambda_i \geq 0$.
        \end{enumerate}
    \end{theorem}
    \centering\notesfig[0.62]{L3-hessian-shapes}
    \\[0.1em]
    \raggedright\small In the eigenbasis, $\by = Q^\top(\bx - \bx^*)$: \ $f(\bx) \approx f(\bx^*) + \frac{1}{2}(\lambda_1y_1^2 + \cdots + \lambda_ny_n^2)$. Some $\lambda_i = 0$: inconclusive.
\end{frame}

\board[2]{Proof of the second-order test}{Taylor: $f(\bx^* + \mathbf{h}) = f(\bx^*) + \frac{1}{2}\mathbf{h}^\top\nabla^2 f(\bx^*)\mathbf{h} + o(\norm{\mathbf{h}}^2)$, and $\mathbf{h}^\top\nabla^2 f(\bx^*)\mathbf{h} \geq \lambda_n\norm{\mathbf{h}}^2$ (Rayleigh, Lecture~6).}

\board{Example: classify the critical points}{$f(x, y) = x^3 + y^3 - 3xy$. Find the critical points and the eigenvalues and eigenvectors of the Hessian there. Is the minimum global?}

\checkslide{$x^3 + y^3 - 3xy$}{
    \sidebyside{0.6}{%
        \small
        \textbf{Critical points.} $\nabla f = (3x^2 - 3y,\ 3y^2 - 3x)^\top = \mathbf{0}$: $y = x^2$ and $x = y^2 = x^4$, so $x(x^3 - 1) = 0$: $x = 0$ or $x = 1$. Critical points $(0, 0)$ and $(1, 1)$.
        \\[0.3em]
        \textbf{Hessian} $\nabla^2 f = \begin{pmatrix} 6x & -3 \\ -3 & 6y \end{pmatrix}$.
        \\[0.3em]
        \textbf{At $(1, 1)$:} $H = \begin{pmatrix} 6 & -3 \\ -3 & 6 \end{pmatrix}$, \ $H(1, 1)^\top = (3, 3)^\top$, $H(1, -1)^\top = (9, -9)^\top$: eigenvalues $3$ and $9$, both $> 0$: strict local minimum, $f(1, 1) = 1 + 1 - 3 = -1$. Not global: $f(x, 0) = x^3 \to -\infty$.
        \\[0.3em]
        \textbf{At $(0, 0)$:} $H = \begin{pmatrix} 0 & -3 \\ -3 & 0 \end{pmatrix}$: eigenvalues $-3$ along $(1, 1)$ and $3$ along $(1, -1)$: a saddle. Indeed $f(t, t) = 2t^3 - 3t^2 \approx -3t^2$ and $f(t, -t) = 3t^2 = \frac{1}{2}\cdot 3 \cdot\norm{(t, -t)}^2$.
    }{%
        \centering\notesfig[1]{L3-critical-points}
    }
}

\begin{frame}{Quadratic functions}
    \begin{example}[][ex:quadratic-min]
        \small $f(\bx) = \frac{1}{2}\bx^\top A\bx - \mathbf{b}^\top\bx$ with $A$ SPD: \ $\nabla f = A\bx - \mathbf{b}$, unique critical point $\bx^* = A^{-1}\mathbf{b}$. Using $\mathbf{b} = A\bx^*$:
        \[f(\bx) - f(\bx^*) = \tfrac{1}{2}\bx^\top A\bx - \bx^{*\top}A\bx - \tfrac{1}{2}\bx^{*\top}A\bx^* + \bx^{*\top}A\bx^* = \tfrac{1}{2}(\bx - \bx^*)^\top A(\bx - \bx^*) > 0\]
        for $\bx \neq \bx^*$: the unique global minimum. Level sets: ellipsoids along the eigenvectors of $A$. \ E.g.\ $A = \diag(1, 10)$, $\mathbf{b} = (1, 10)^\top$: $\bx^* = (1, 1)^\top$, $f - f^* = \frac{1}{2}\big((x_1 - 1)^2 + 10(x_2 - 1)^2\big)$.
    \end{example}
    \small
    \begin{itemize}
        \item Solving $A\bx = \mathbf{b}$ with SPD $A$ $=$ minimizing a quadratic: the start of iterative solvers (conjugate gradients).
        \item Least squares $\norm{\by - \Phi\boldsymbol{\theta}}^2$ is of this form with $A = 2\Phi^\top\Phi$, $\mathbf{b} = 2\Phi^\top\by$.
        \item Every smooth function looks like this near a strict local minimum, with $A = \nabla^2 f(\bx^*)$: what we learn about quadratics describes the end of every optimization.
    \end{itemize}
\end{frame}

% ------------------------------------------------------------------ Gradient descent
\topic{\secref{sec:gd}}{Gradient descent and Newton's method}

\begin{frame}{Gradient descent, derived for the mean squared error}
    \begin{definition}[Gradient descent][def:gd]
        \small $\bx_{k+1} = \bx_k - \eta\,\nabla f(\bx_k)$, with step size (learning rate) $\eta > 0$.
    \end{definition}
    \begin{example}[Linear regression][ex:gd-mse]
        \small $L(\bw) = \frac{1}{N}\norm{X\bw - \by}^2$, \ $\nabla L(\bw) = \frac{2}{N}X^\top(X\bw - \by)$:
        \[\bw_{k+1} = \bw_k - \eta\,\frac{2}{N}\sum_{i=1}^N(\bw_k^\top\bx_i - y_i)\,\bx_i.\]
        Each data point pulls $\bw$ along its input, in proportion to its residual. One step costs $O(Nd)$: no $d \times d$ system to solve.
    \end{example}
\end{frame}

\board{Example: gradient descent by hand}{Model $y = wx$, data $(1, 2)$, $(2, 3)$, $(3, 5)$. Find $L'(w)$, the minimum $w^*$, three steps with $\eta = 0.1$ from $w_0 = 0$, and the step sizes for which GD converges.}

\checkslide{gradient descent for $y = wx$}{
    \small
    $L(w) = \frac{1}{3}\sum_i(wx_i - y_i)^2$, \ $\sum_ix_i^2 = 1 + 4 + 9 = 14$, \ $\sum_ix_iy_i = 2 + 6 + 15 = 23$:
    \[L'(w) = \tfrac{2}{3}\textstyle\sum_i(wx_i - y_i)x_i = \tfrac{2}{3}(14w - 23), \qquad L'' = \tfrac{28}{3}, \qquad w^* = \tfrac{23}{14} \approx 1.6429.\]
    \textbf{Three steps}, $\eta = 0.1$, $w_0 = 0$:
    \[\begin{aligned}
        w_1 &= 0 - 0.1 \cdot \tfrac{2}{3}(0 - 23) = 1.5333, \\
        w_2 &= 1.5333 - 0.1 \cdot \tfrac{2}{3}(21.4667 - 23) = 1.5333 + 0.1022 = 1.6356, \\
        w_3 &= 1.6356 - 0.1 \cdot \tfrac{2}{3}(22.8978 - 23) = 1.6356 + 0.0068 = 1.6424.
    \end{aligned}\]
    \textbf{Errors} $w_k - w^*$: $-1.643,\ -0.110,\ -0.0073,\ -0.0005$: each is $\frac{1}{15}$ of the previous one, since
    \[w_{k+1} - w^* = (1 - \eta L'')(w_k - w^*), \qquad 1 - 0.1 \cdot \tfrac{28}{3} = \tfrac{1}{15}.\]
    Converges $\iff |1 - \eta L''| < 1 \iff 0 < \eta < \frac{2}{28/3} = \frac{3}{14} \approx 0.21$; fastest at $\eta = \frac{3}{28}$ (one step).
}

\begin{frame}{Smooth functions: gradient descent always makes progress}
    \small
    $f$ is \textbf{$L$-smooth} if $\norm{\nabla f(\bx) - \nabla f(\by)} \leq L\norm{\bx - \by}$ (for $C^2$: Hessian eigenvalues in $[-L, L]$).
    \begin{lemma}[Descent lemma][lem:descent]
        \small $f(\by) \leq f(\bx) + \nabla f(\bx)^\top(\by - \bx) + \frac{L}{2}\norm{\by - \bx}^2$.
    \end{lemma}
    \begin{theorem}[][thm:gd-smooth]
        \small If $f$ is $L$-smooth, $f \geq f^*$, and $0 < \eta \leq \frac{1}{L}$, then
        \[f(\bx_{k+1}) \leq f(\bx_k) - \frac{\eta}{2}\norm{\nabla f(\bx_k)}^2, \qquad \min_{k < K}\norm{\nabla f(\bx_k)}^2 \leq \frac{2(f(\bx_0) - f^*)}{\eta K}.\]
    \end{theorem}
    GD finds points with small gradient, but not necessarily a minimum (saddles, bad local minima). For neural networks, this is essentially the best general guarantee.
\end{frame}

\board[2]{Proof of the descent lemma and the theorem}{Descent lemma: $f(\by) - f(\bx) = \int_0^1\nabla f(\bx + t(\by - \bx))^\top(\by - \bx)\,dt$ and Cauchy--Schwarz. \quad Theorem: put $\by = \bx_k - \eta\nabla f(\bx_k)$; then sum over $k$ (telescoping).}

\begin{frame}{Backtracking line search}
    \small
    $L$ is unknown in practice. \textbf{Backtracking:} start with $\eta = \bar\eta$ and halve it until the Armijo condition
    \[f\big(\bx_k - \eta\nabla f(\bx_k)\big) \leq f(\bx_k) - \frac{\eta}{2}\norm{\nabla f(\bx_k)}^2\]
    holds; then step.
    \begin{proposition}[][prop:backtracking]
        \small For $L$-smooth $f \geq f^*$, backtracking stops after finitely many halvings with $\eta_k \geq \eta_{\min} = \min(\bar\eta, \frac{1}{2L})$, and $\min_{k < K}\norm{\nabla f(\bx_k)}^2 \leq \frac{2(f(\bx_0) - f^*)}{\eta_{\min}K}$.
    \end{proposition}
    Why it stops: by the theorem, every $\eta \leq \frac{1}{L}$ satisfies the condition, and halving reaches $(\frac{1}{2L}, \frac{1}{L}]$.
\end{frame}

\begin{frame}{Gradient descent on quadratic functions}
    \sidebyside{0.55}{%
        \small
        $f = \frac{1}{2}\bx^\top A\bx - \mathbf{b}^\top\bx$, $\mathbf{e}_k = \bx_k - \bx^*$: \ $\mathbf{e}_{k+1} = (I - \eta A)\mathbf{e}_k$.
        \begin{theorem}[][thm:gd-quadratic]
            \small With $\kappa = \lambda_{\max}/\lambda_{\min}$:
            \begin{enumerate}
                \item $\norm{\mathbf{e}_k} \leq \rho^k\norm{\mathbf{e}_0}$, $\rho = \max_i|1 - \eta\lambda_i|$;
                \item convergence for all $\bx_0$ $\iff 0 < \eta < \frac{2}{\lambda_{\max}}$;
                \item best $\eta = \frac{2}{\lambda_{\max} + \lambda_{\min}}$, \ $\rho = \frac{\kappa - 1}{\kappa + 1}$.
            \end{enumerate}
        \end{theorem}
    }{%
        \centering\notesfig[1]{L3-gd}
    }
    \begin{example}[][ex:gd-quadratic]
        \small $\frac{1}{2}(x_1^2 + 10x_2^2)$: $\eta < 0.2$; best $\eta = \frac{2}{11}$, $\rho = \frac{9}{11}$; $69$ steps to gain a factor $10^{-6}$. For $A = I$: one step.
    \end{example}
\end{frame}

\board{Proof for quadratic functions}{$\mathbf{e}_{k+1} = (I - \eta A)\mathbf{e}_k$; in the eigenbasis of $A$ each coordinate is multiplied by $1 - \eta\lambda_i$. Minimize $\max(|1 - \eta\lambda_{\max}|, |1 - \eta\lambda_{\min}|)$ over $\eta$.}

\checkslide{gradient descent on $\frac{1}{2}(x_1^2 + 10x_2^2)$}{
    \small
    $A = \diag(1, 10)$, $\bx^* = \mathbf{0}$: \ $\bx_{k+1} = (I - \eta A)\bx_k$, i.e., $x_1 \mapsto (1 - \eta)x_1$, \ $x_2 \mapsto (1 - 10\eta)x_2$.
    \begin{center}
        \renewcommand{\arraystretch}{1.25}
        \begin{tabular}{@{}lccl@{}}
            \toprule
            $\eta$ & factor on $x_1$ & factor on $x_2$ & \\
            \midrule
            $0.05$ & $0.95$ & $0.5$ & slow along the valley \\
            $\frac{2}{11} \approx 0.18$ & $\frac{9}{11}$ & $-\frac{9}{11}$ & best: $\rho = \frac{\kappa - 1}{\kappa + 1} = \frac{9}{11}$ \\
            $0.21$ & $0.79$ & $-1.1$ & diverges: $\eta > \frac{2}{\lambda_{\max}} = 0.2$ \\
            \bottomrule
        \end{tabular}
    \end{center}
    \textbf{Best $\eta$:} $|1 - \eta\lambda_{\min}| = |1 - \eta\lambda_{\max}|$ $\Rightarrow$ $1 - \eta = 10\eta - 1$, $\eta = \frac{2}{11}$. Steps to gain $10^{-6}$: $k \geq \frac{\ln 10^6}{\ln(11/9)} \approx 68.8$, so $69$.
    \\[0.4em]
    \textbf{$\kappa = 25$} (the plot on the next slide): GD $\rho = \frac{24}{26} \approx 0.923$, about $173$ steps for $10^{-6}$; momentum $\rho \approx \frac{\sqrt{25} - 1}{\sqrt{25} + 1} = \frac{2}{3}$, about $35$ steps. \ For $\kappa = 10^4$: about $70\,000$ vs $700$.
}

\begin{frame}{Making gradient descent faster}
    \small
    \textbf{Feature scaling:} standardized columns of $\Phi$ often make $\kappa(\Phi^\top\Phi)$ much smaller. \\
    \textbf{Momentum} (heavy ball): $\bx_{k+1} = \bx_k - \eta\nabla f(\bx_k) + \beta(\bx_k - \bx_{k-1})$: about $\sqrt{\kappa}$ instead of $\kappa$ steps. \\
    \textbf{Preconditioning:} the gradient for $\ip{\bx}{\by}_P$ is $P^{-1}\nabla f$ (Lecture~7); Adam: adaptive diagonal $P$; Newton: $P = \nabla^2 f$.
    \\[0.2em]
    \centering\fig[0.74]{s08_momentum}
\end{frame}

\begin{frame}{Stochastic gradient descent}
    \begin{remark}[][rem:sgd]
        \small $L(\boldsymbol{\theta}) = \frac{1}{N}\sum_{i=1}^N\ell_i(\boldsymbol{\theta})$. Instead of the full gradient (a pass over all data), use a random mini-batch $B$:
        \[\mathbf{g} = \frac{1}{|B|}\sum_{i \in B}\nabla\ell_i(\boldsymbol{\theta}), \qquad \E[\mathbf{g}] = \nabla L(\boldsymbol{\theta}).\]
    \end{remark}
    \small
    \begin{itemize}
        \item On average SGD moves in the right direction; the noise prevents it from settling at the minimum unless $\eta$ decreases.
        \item Lecture~10 (self-study \S4.7.3): convergence for strongly convex $f$ with $\eta_k = \frac{1}{\mu(k + 1)}$ at the rate $\frac{1}{k}$; with constant $\eta$, a neighborhood of size $\propto \eta$.
        \item SGD with momentum and preconditioning (Adam) trains almost all neural networks.
    \end{itemize}
\end{frame}

\begin{frame}{Newton's method}
    \small
    Minimize the quadratic model $m_k(\mathbf{h}) = f(\bx_k) + \nabla f(\bx_k)^\top\mathbf{h} + \frac{1}{2}\mathbf{h}^\top\nabla^2 f(\bx_k)\mathbf{h}$:
    \[\bx_{k+1} = \bx_k - \nabla^2 f(\bx_k)^{-1}\nabla f(\bx_k) \qquad \text{(solve } H_k\mathbf{h} = -\nabla f(\bx_k)\text{, do not invert)}.\]
    \begin{theorem}[Quadratic convergence][thm:newton]
        \small If $\nabla^2 f(\bx^*) \succeq \mu I$, $\mu > 0$, the Hessian is $L$-Lipschitz near $\bx^*$, and $\bx_0$ is close enough:
        \(\norm{\bx_{k+1} - \bx^*} \leq \frac{L}{\mu}\norm{\bx_k - \bx^*}^2.\)
    \end{theorem}
    \begin{example}[][ex:newton]
        \small $f(x) = e^x - 2x$, $x^* = \log 2$, $x_{k+1} = x_k - 1 + 2e^{-x_k}$: errors $0.69,\ 0.31,\ 4.3 \cdot 10^{-2},\ 9.0 \cdot 10^{-4},\ 4.0 \cdot 10^{-7},\ 8.0 \cdot 10^{-14}$: the number of correct digits doubles.
    \end{example}
    Drawbacks: $n^2$ entries and $n^3$ work; far from a minimum it may go to a saddle or a maximum. Newton $=$ gradient descent for the inner product of the Hessian.
\end{frame}

\board{Derive the Newton step}{Minimize $m_k(\mathbf{h})$: set $\nabla m_k(\mathbf{h}) = \mathbf{0}$. Why does Newton solve a quadratic in one step, for every $\kappa$? \quad One variable: $x_{k+1} = x_k - f'(x_k)/f''(x_k)$.}

\checkslide{Newton and Gauss--Newton with numbers}{
    \small
    \textbf{Derivation.} $\nabla m_k(\mathbf{h}) = \nabla f(\bx_k) + H_k\mathbf{h} = \mathbf{0}$ $\Rightarrow$ $\mathbf{h} = -H_k^{-1}\nabla f(\bx_k)$. For a quadratic, $m_0 = f$, so one step lands on $\bx^*$.
    \\[0.3em]
    \textbf{Newton for $f(x) = e^x - 2x$:} $f' = e^x - 2$, $f'' = e^x$, \ $x_{k+1} = x_k - \frac{e^{x_k} - 2}{e^{x_k}} = x_k - 1 + 2e^{-x_k}$:
    \[x_1 = 0 - 1 + 2 = 1, \quad x_2 = 1 - 1 + 2e^{-1} \approx 0.735759, \quad x_3 \approx 0.735759 - 1 + 2 \cdot 0.479142 \approx 0.694042,\]
    then $0.6931476$ and $0.693147180560$; errors $0.69, 0.31, 4.3 \cdot 10^{-2}, 9.0 \cdot 10^{-4}, 4.0 \cdot 10^{-7}, 8 \cdot 10^{-14}$: each about the square of the previous one (times a constant).
    \\[0.4em]
    \textbf{Gauss--Newton} for $y = e^{\theta x}$, data $(0, 1), (1, 2), (2, 4)$, $\theta_0 = 1$: residuals $r_i = e^{\theta x_i} - y_i$, $J_i = x_ie^{\theta x_i}$:
    \[\mathbf{r} \approx (0, 0.718, 3.389), \quad J \approx (0, 2.718, 14.778), \quad h = -\frac{\sum_iJ_ir_i}{\sum_iJ_i^2} \approx -\frac{52.04}{225.78} \approx -0.2305,\]
    so $\theta_1 \approx 0.7695$; then $0.69856$, $0.693175$, \ldots $\to \log 2$: quadratic, since the residuals vanish at $\theta^*$.
}

\begin{frame}{Newton-type methods in practice}
    \small
    \textbf{Nonlinear least squares} $f(\boldsymbol{\theta}) = \frac{1}{2}\norm{\mathbf{r}(\boldsymbol{\theta})}^2$, $J = D\mathbf{r}$:
    \[\nabla f = J^\top\mathbf{r}, \qquad \nabla^2 f = J^\top J + \textstyle\sum_i r_i\nabla^2 r_i.\]
    \begin{itemize}
        \item \textbf{Gauss--Newton} drops the second term: solve $J^\top J\mathbf{h} = -J^\top\mathbf{r}$ (a least squares problem). Always a descent direction; for linear models, one step.
        \item \textbf{Levenberg--Marquardt}: $(J^\top J + \lambda I)\mathbf{h} = -J^\top\mathbf{r}$, between Gauss--Newton and gradient descent.
        \item Example: fit $y = e^{\theta x}$ to $(0, 1), (1, 2), (2, 4)$ from $\theta_0 = 1$: errors $0.31,\ 0.076,\ 5.4 \cdot 10^{-3},\ 2.8 \cdot 10^{-5},\ 7.8 \cdot 10^{-10}$.
    \end{itemize}
    \textbf{Hessian-vector products} $\nabla^2 f(\bx)\bv = \frac{d}{dt}\nabla f(\bx + t\bv)|_{t=0}$ cost a few gradients with AD (forward over reverse): conjugate gradients, the top eigenvalue of a network's Hessian by power iteration. Quasi-Newton (BFGS, L-BFGS) builds the Hessian from gradients.
\end{frame}

% ------------------------------------------------------------------ Convexity
\topic{\secref{sec:convexity}}{Convexity}

\begin{frame}{Convex sets and functions}
    \begin{definition}[][def:convex]
        \small $C$ is convex if it contains the segment between any two of its points. $f$ is convex if
        \[f(\theta\bx + (1 - \theta)\by) \leq \theta f(\bx) + (1 - \theta)f(\by), \qquad \theta \in [0, 1].\]
    \end{definition}
    \sidebyside{0.45}{%
        \small
        Examples: affine functions; every norm; $x^2$, $|x|$, $e^x$, $-\log x$. Not convex: $x^3$.
        \begin{theorem}[First order][thm:convex-first-order]
            \small Differentiable $f$ is convex $\iff$
            \(f(\by) \geq f(\bx) + \nabla f(\bx)^\top(\by - \bx).\)
        \end{theorem}
    }{%
        \centering\notesfig[1]{L3-convex}
    }
    \small A convex function lies below its chords and above its tangent planes.
\end{frame}

\begin{frame}{Why convexity matters}
    \begin{corollary}[][cor:convex-global]
        \small For convex differentiable $f$: \ $\nabla f(\bx^*) = \mathbf{0}$ $\Rightarrow$ $\bx^*$ is a \emph{global} minimum.
    \end{corollary}
    \begin{proposition}[][prop:convex-minima]
        \small Every local minimum of a convex function is global; the set of minima is convex; strictly convex $\Rightarrow$ at most one minimum.
    \end{proposition}
    \begin{theorem}[Second order][thm:convex-second-order]
        \small $C^2$ function $f$ is convex $\iff$ $\nabla^2 f(\bx) \succeq 0$ for all $\bx$; \ $\nabla^2 f \succ 0$ everywhere $\Rightarrow$ strictly convex.
    \end{theorem}
    \begin{proposition}[][prop:convex-operations]
        \small $\alpha f + \beta g$ ($\alpha, \beta \geq 0$), \ $f(A\bx + \mathbf{b})$, \ $\max(f, g)$ are convex if $f$, $g$ are.
    \end{proposition}
\end{frame}

\board{Proofs: tangent plane $\Rightarrow$ global minimum; local $=$ global}{(1) Put $\nabla f(\bx^*) = \mathbf{0}$ into the first-order characterization. \quad (2) If $f(\by) < f(\bx^*)$, look at $f$ on the segment from $\bx^*$ to $\by$.}

\begin{frame}{Convex losses, and non-convex networks}
    \begin{example}[Convex losses][ex:convex-losses]
        \small
        \begin{itemize}
            \item least squares: Hessian $2\Phi^\top\Phi \succeq 0$; with a ridge penalty $\lambda\norm{\boldsymbol{\theta}}^2$: $\succ 0$, a unique minimum;
            \item logistic regression: $\frac{1}{N}\sum_ip_i(1 - p_i)\bx_i\bx_i^\top \succeq 0$;
            \item log-sum-exp: $\bv^\top(\diag(\mathbf{p}) - \mathbf{p}\mathbf{p}^\top)\bv = \sum_ip_iv_i^2 - (\sum_ip_iv_i)^2 \geq 0$ (Cauchy--Schwarz).
        \end{itemize}
    \end{example}
    \begin{remark}[Neural networks are not convex]
        \small Swapping two hidden units gives the same function, hence the same loss. If the loss were convex, the midpoint of the two minima would be a minimum too, but there the two units are identical: effectively one unit less. Yet gradient methods work remarkably well in practice.
    \end{remark}
\end{frame}

\begin{frame}{Strong convexity and convergence rates}
    \begin{definition}[][def:strongly-convex]
        \small $f$ is $\mu$-strongly convex if $f(\by) \geq f(\bx) + \nabla f(\bx)^\top(\by - \bx) + \frac{\mu}{2}\norm{\by - \bx}^2$ \ (for $C^2$: $\nabla^2 f \succeq \mu I$).
    \end{definition}
    \begin{lemma}[][lem:pl]
        \small $\nabla f(\bx)^\top(\bx - \bx^*) \geq \mu\norm{\bx - \bx^*}^2$ \quad and \quad $f(\bx) - f^* \leq \frac{1}{2\mu}\norm{\nabla f(\bx)}^2$ \ (Polyak--{\L}ojasiewicz).
    \end{lemma}
    \begin{theorem}[Rates of gradient descent][thm:gd-rates]
        \small $f$ convex and $L$-smooth, $\eta = \frac{1}{L}$:
        \begin{enumerate}
            \item $f(\bx_k) - f^* \leq \dfrac{L\norm{\bx_0 - \bx^*}^2}{2k}$;
            \item if also $\mu$-strongly convex: $f(\bx_k) - f^* \leq \big(1 - \frac{\mu}{L}\big)^k(f(\bx_0) - f^*)$.
        \end{enumerate}
    \end{theorem}
    \small $L/\mu$ plays the role of the condition number $\kappa$.
\end{frame}

\board[2]{Proof of the linear rate}{Descent lemma: $f(\bx_{k+1}) \leq f(\bx_k) - \frac{1}{2L}\norm{\nabla f(\bx_k)}^2$. \ PL: $\norm{\nabla f(\bx_k)}^2 \geq 2\mu(f(\bx_k) - f^*)$. Combine. \quad PL itself: minimize the right side of the strong convexity inequality over $\by$.}

\checkslide{the linear rate}{
    \small
    \textbf{PL.} Minimize the right side of $f(\by) \geq f(\bx) + \nabla f(\bx)^\top(\by - \bx) + \frac{\mu}{2}\norm{\by - \bx}^2$ over $\by$: at $\by - \bx = -\frac{1}{\mu}\nabla f(\bx)$,
    \[f^* \geq f(\bx) - \tfrac{1}{\mu}\norm{\nabla f}^2 + \tfrac{1}{2\mu}\norm{\nabla f}^2 = f(\bx) - \tfrac{1}{2\mu}\norm{\nabla f(\bx)}^2.\]
    \textbf{Combine} (descent lemma with $\eta = \frac{1}{L}$, then PL):
    \[f(\bx_{k+1}) - f^* \leq f(\bx_k) - f^* - \tfrac{1}{2L}\norm{\nabla f(\bx_k)}^2 \leq \big(1 - \tfrac{\mu}{L}\big)\big(f(\bx_k) - f^*\big).\]
    \textbf{Example.} $f = \frac{1}{2}(x_1^2 + 10x_2^2)$: Hessian eigenvalues $1$ and $10$, so $\mu = 1$, $L = 10$, $\eta = 0.1$: the gap shrinks at least by $0.9$ per step; $10^{-6}$ after $\frac{\ln 10^6}{-\ln 0.9} \approx 131$ steps. (With the better step $\frac{2}{11}$ the error itself shrinks by $\frac{9}{11}$, the gap by $(\frac{9}{11})^2 \approx 0.67$.)
    \\[0.3em]
    \textbf{Without strong convexity}, e.g.\ $f(x) = x^4$: $f''(0) = 0$, and the gap only decays like $\frac{1}{k}$ (at least, by the $O(1/k)$ theorem).
}

% ------------------------------------------------------------------ Prox
\topic{\secref{sec:prox}}{Nonsmooth convex functions and the lasso}

\begin{frame}{Subgradients}
    \begin{definition}[Subgradient][def:subgradient]
        \small $\mathbf{g} \in \partial f(\bx)$ if $f(\by) \geq f(\bx) + \mathbf{g}^\top(\by - \bx)$ for all $\by$.
    \end{definition}
    \begin{proposition}[][prop:subgradients]
        \small Differentiable: $\partial f(\bx) = \{\nabla f(\bx)\}$. \ $\partial|x| = \{\sign x\}$ for $x \neq 0$, $\partial|x|(0) = [-1, 1]$. \ For $\norm{\bx}_1$: $g_i = \sign(x_i)$ if $x_i \neq 0$, $g_i \in [-1, 1]$ if $x_i = 0$.
    \end{proposition}
    \begin{theorem}[Optimality][thm:composite-optimality]
        \small $f$ convex differentiable, $h$ convex: \ $\bx^*$ minimizes $f + h$ $\iff$ $-\nabla f(\bx^*) \in \partial h(\bx^*)$.
    \end{theorem}
    \small
    For $h = \lambda\norm{\cdot}_1$: $w_i^* = 0$ is optimal as long as $|\partial_i f(\bw^*)| \leq \lambda$. A whole interval of gradients is allowed at $0$, so many weights are exactly $0$.
\end{frame}

\begin{frame}{The proximal gradient method}
    \small
    \begin{definition}[Proximal operator][def:prox]
        $\operatorname{prox}_{\eta h}(\bv) = \argmin_{\bu}\big(h(\bu) + \frac{1}{2\eta}\norm{\bu - \bv}^2\big)$.
    \end{definition}
    \begin{proposition}[Soft thresholding][prop:soft-threshold]
        $\operatorname{prox}_{\eta\lambda\norm{\cdot}_1}(\bv)_i = S_{\eta\lambda}(v_i)$, \ $S_\tau(v) = \sign(v)\max(|v| - \tau, 0)$.
    \end{proposition}
    \begin{theorem}[Proximal gradient (ISTA)][thm:prox-gradient]
        $\bx_{k+1} = \operatorname{prox}_{\eta h}(\bx_k - \eta\nabla f(\bx_k))$, $\eta \leq \frac{1}{L}$: \ $\bx^*$ is optimal iff it is a fixed point, and $F = f + h$ decreases: $F(\bx_{k+1}) \leq F(\bx_k) - \frac{1}{2\eta}\norm{\bx_{k+1} - \bx_k}^2$.
    \end{theorem}
    Gradient step on the smooth part, prox step on the nonsmooth part. With $h = 0$: gradient descent; with $h = $ indicator of a set: projected gradient descent.
\end{frame}

\board{Derive soft thresholding}{Minimize $\lambda|u| + \frac{1}{2}(u - v)^2$ over $u \in \R$: use the optimality condition $v - u \in \lambda\,\partial|u|$; three cases.}

\checkslide{soft thresholding}{
    \small
    Optimality ($f = \frac{1}{2}(u - v)^2$, $h = \lambda|u|$): \ $-f'(u) = v - u \in \lambda\,\partial|u|$.
    \begin{itemize}
        \item $u > 0$: $\partial|u| = \{1\}$, so $u = v - \lambda$; consistent iff $v > \lambda$.
        \item $u < 0$: $\partial|u| = \{-1\}$, so $u = v + \lambda$; consistent iff $v < -\lambda$.
        \item $u = 0$: $\partial|0| = [-1, 1]$, so $v \in [-\lambda, \lambda]$.
    \end{itemize}
    Exactly one case applies for each $v$: \ $S_\lambda(v) = \sign(v)\max(|v| - \lambda, 0)$.
    \\[0.4em]
    \textbf{Numbers} ($\lambda = 1$, orthonormal features, least squares weights $(3, -0.5, 1.2, 0.1)$):
    \begin{center}
        \renewcommand{\arraystretch}{1.2}
        \begin{tabular}{@{}lcccc@{}}
            \toprule
            least squares $v$ & $3$ & $-0.5$ & $1.2$ & $0.1$ \\
            lasso $S_1(v)$ & $3 - 1 = 2$ & $0$ ($|v| \leq 1$) & $1.2 - 1 = 0.2$ & $0$ \\
            ridge $\frac{v}{1 + 2\lambda} = \frac{v}{3}$ & $1$ & $-0.167$ & $0.4$ & $0.033$ \\
            \bottomrule
        \end{tabular}
    \end{center}
    The lasso keeps $2$ of $4$ features; ridge keeps all, just smaller. ISTA step: $\bw \leftarrow S_{\eta\lambda}(\bw - \eta\nabla f(\bw))$.
}

\begin{frame}{Why the lasso gives sparse weights}
    \centering\fig[0.64]{s08_lasso}
    \\[0.1em]
    \raggedright
    \begin{example}[][ex:lasso]
        \small If $X^\top X = I$: lasso weights $w_i^* = S_\lambda\big((X^\top\by)_i\big)$; $(3, -0.5, 1.2, 0.1)$ with $\lambda = 1$ becomes $(2, 0, 0.2, 0)$. Ridge: $\frac{1}{1 + 2\lambda}X^\top\by$, all shrunk, none $0$.
    \end{example}
    \small \textcolor{mlcolor}{Penalty $\lambda\norm{\bw}$ $\leftrightarrow$ constraint $\norm{\bw} \leq t$ (KKT, next part): the loss's level sets first touch the $\ell_1$ ball at a corner.}
\end{frame}

% ------------------------------------------------------------------ Constrained
\topic{\secref{sec:constrained}}{Constrained optimization}

\begin{frame}{Lagrange multipliers}
    \sidebyside{0.56}{%
        \begin{theorem}[Lagrange][thm:lagrange]
            \small Let $\bx^*$ be a local minimum of $f$ on $\{g_1 = \cdots = g_m = 0\}$, with $\nabla g_1(\bx^*), \ldots, \nabla g_m(\bx^*)$ linearly independent. Then there are unique $\lambda_i$ with
            \[\nabla f(\bx^*) + \sum_{i=1}^m\lambda_i\nabla g_i(\bx^*) = \mathbf{0}.\]
        \end{theorem}
        \small
        $\nabla f$ is orthogonal to the constraint set: the level set of $f$ touches it.
    }{%
        \centering\notesfig[0.9]{L3-lagrange}
    }
    \begin{example}[][ex:lagrange-circle]
        \small $\max x + y$ on $x^2 + y^2 = 1$: $(1, 1) + \lambda(2x, 2y) = \mathbf{0}$, so $x = y = \pm\frac{1}{\sqrt2}$; the maximum is $\sqrt{2}$ (as Cauchy--Schwarz predicts).
    \end{example}
\end{frame}

\board[2]{Two Lagrange problems from linear algebra}{(1) Maximize $\bx^\top A\bx$ subject to $\bx^\top\bx = 1$ ($A$ symmetric): the candidates are eigenvectors, the multiplier is the eigenvalue (PCA). \\ (2) Minimize $\norm{\bx}^2$ subject to $A\bx = \mathbf{b}$ (independent rows): $\bx^* = A^\top(AA^\top)^{-1}\mathbf{b}$, in $\row(A)$.}

\checkslide{Lagrange multipliers with numbers}{
    \small
    \textbf{(1) Eigenvalues.} $\nabla(\bx^\top A\bx) - \lambda\nabla(\bx^\top\bx) = 2A\bx - 2\lambda\bx = \mathbf{0}$: $A\bx = \lambda\bx$, and $f = \bx^\top A\bx = \lambda$ there. For $A = \begin{pmatrix} 2.5 & 1.5 \\ 1.5 & 2.5 \end{pmatrix}$: candidates $\pm\frac{1}{\sqrt2}(1, 1)^\top$ ($f = 4$, max) and $\pm\frac{1}{\sqrt2}(1, -1)^\top$ ($f = 1$, min).
    \\[0.5em]
    \textbf{(2) Minimum norm.} $2\bx + A^\top\boldsymbol{\lambda} = \mathbf{0}$, so $\bx = A^\top\boldsymbol{\nu}$; then $AA^\top\boldsymbol{\nu} = \mathbf{b}$. Example: $x_1 + x_2 = 1$, $x_2 + x_3 = 1$:
    \[A = \begin{pmatrix} 1 & 1 & 0 \\ 0 & 1 & 1 \end{pmatrix}, \quad AA^\top = \begin{pmatrix} 2 & 1 \\ 1 & 2 \end{pmatrix}, \quad \boldsymbol{\nu} = (AA^\top)^{-1}\begin{pmatrix} 1 \\ 1 \end{pmatrix} = \tfrac{1}{3}\begin{pmatrix} 2 & -1 \\ -1 & 2 \end{pmatrix}\begin{pmatrix} 1 \\ 1 \end{pmatrix} = \begin{pmatrix} 1/3 \\ 1/3 \end{pmatrix},\]
    \[\bx^* = A^\top\boldsymbol{\nu} = \big(\tfrac{1}{3}, \tfrac{2}{3}, \tfrac{1}{3}\big)^\top: \quad \tfrac{1}{3} + \tfrac{2}{3} = 1 \ \checkmark, \quad \norm{\bx^*}^2 = \tfrac{6}{9} = \tfrac{2}{3}.\]
    Every other solution is $\bx^* + t(1, -1, 1)^\top$ (the null space), with norm$^2$ $\frac{2}{3} + 3t^2$: larger.
}

\begin{frame}{Inequality constraints: the KKT conditions}
    \begin{theorem}[Karush--Kuhn--Tucker][thm:kkt]
        \small Let $\bx^*$ be a local minimum of $f$ subject to $g_i(\bx) \leq 0$, with independent gradients of the active constraints. Then there are $\mu_i$ with
        \begin{enumerate}
            \item stationarity: $\nabla f(\bx^*) + \sum_i\mu_i\nabla g_i(\bx^*) = \mathbf{0}$;
            \item primal feasibility: $g_i(\bx^*) \leq 0$; \quad dual feasibility: $\mu_i \geq 0$;
            \item complementary slackness: $\mu_ig_i(\bx^*) = 0$ (inactive constraints have $\mu_i = 0$).
        \end{enumerate}
        If $f$ and all $g_i$ are convex, every KKT point is a global minimum.
    \end{theorem}
    \small $\mu_i \geq 0$: the constraint pushes back, it cannot pull. $\mu_i$ is the price of the constraint: relaxing $g_i \leq 0$ to $g_i \leq \varepsilon$ lowers the minimum by about $\mu_i\varepsilon$.
\end{frame}

\board[2]{Examples of KKT}{(1) Minimize $(x - 2)^2$ subject to $x \leq 1$. What if the constraint is $x \leq 3$? \\ (2) Project $\mathbf{p}$ with $\norm{\mathbf{p}} > 1$ onto the unit ball: minimize $\norm{\bx - \mathbf{p}}^2$ subject to $\norm{\bx}^2 - 1 \leq 0$. Numbers: $\mathbf{p} = (3, 4)^\top$.}

\checkslide{KKT examples}{
    \small
    (1) $2(x - 2) + \mu = 0$, $x \leq 1$, $\mu \geq 0$, $\mu(x - 1) = 0$. \ $\mu = 0$ gives $x = 2$, infeasible; so $x^* = 1$, $\mu = 2$. Relaxing to $x \leq 1 + \varepsilon$: minimum $1 - 2\varepsilon + \varepsilon^2$, rate $\mu = 2$. For $x \leq 3$: $x^* = 2$, inactive, $\mu = 0$.
    \\[0.6em]
    (2) $2(\bx - \mathbf{p}) + 2\mu\bx = \mathbf{0}$, so $\bx = \frac{\mathbf{p}}{1 + \mu}$; $\mu = 0$ is infeasible, so $\norm{\bx} = 1$, $\mu = \norm{\mathbf{p}} - 1$:
    \[\bx^* = \frac{\mathbf{p}}{\norm{\mathbf{p}}}; \qquad \mathbf{p} = (3, 4)^\top: \ \norm{\mathbf{p}} = 5, \ \bx^* = (0.6, 0.8)^\top, \ \mu = 4.\]
    Check stationarity: $2(\bx^* - \mathbf{p}) + 2\mu\bx^* = 2(-2.4, -3.2) + 8(0.6, 0.8) = (0, 0)$ \checkmark; $\mu = 4 \geq 0$, constraint active.
    \textcolor{mlcolor}{Projected gradient descent alternates gradient steps with such projections (e.g., weight constraints, trust regions).}
}

\begin{frame}{Support vector machines}
    \sidebyside{0.56}{%
        \small
        Separable data $(\bx_i, y_i)$, $y_i = \pm 1$. Margin of $\bw^\top\bx + b = 0$ is $\frac{1}{\norm{\bw}}$ (Lecture~6) if we normalize $\min_iy_i(\bw^\top\bx_i + b) = 1$:
        \[\min\ \tfrac{1}{2}\norm{\bw}^2 \quad \text{s.t.} \quad 1 - y_i(\bw^\top\bx_i + b) \leq 0.\]
        KKT: \ $\bw = \sum_i\mu_iy_i\bx_i$, and $\mu_i > 0$ only if $y_i(\bw^\top\bx_i + b) = 1$: the points on the margin, the \textbf{support vectors}. The classifier depends only on them.
        \\[0.3em]
        In the picture: $\bw = (1, 1)^\top$, $b = -2$. Support vectors $(1.5, 1.5)$: $1.5 + 1.5 - 2 = 1$; \ $(0.5, 0.5)$, $y = -1$: $-(0.5 + 0.5 - 2) = 1$. Margin $\frac{1}{\norm{\bw}} = \frac{1}{\sqrt2} \approx 0.71$.
    }{%
        \centering\fig[0.95]{s08_svm}
    }
\end{frame}

\begin{frame}{Lagrangian duality}
    \small
    Lagrangian $\mathcal{L}(\bx, \boldsymbol{\mu}) = f(\bx) + \sum_i\mu_ig_i(\bx)$; \ dual function $d(\boldsymbol{\mu}) = \inf_{\bx}\mathcal{L}(\bx, \boldsymbol{\mu})$ for $\boldsymbol{\mu} \geq \mathbf{0}$: always concave.
    \begin{theorem}[Weak and strong duality][thm:duality]
        \begin{enumerate}
            \item $d(\boldsymbol{\mu}) \leq f(\bx)$ for every feasible $\bx$ and $\boldsymbol{\mu} \geq \mathbf{0}$; hence $d^* \leq p^*$.
            \item If $f$, $g_i$ are convex and $\bx^*$, $\boldsymbol{\mu}^*$ satisfy KKT, then they solve the primal and the dual problem, and $d^* = p^*$.
        \end{enumerate}
    \end{theorem}
    \begin{example}[][ex:dual-1d]
        $\min(x - 2)^2$ s.t. $x \leq 1$: $d(\mu) = \mu - \frac{\mu^2}{4}$, maximal at $\mu = 2$ with $d^* = 1 = p^*$.
    \end{example}
    Every dual feasible $\boldsymbol{\mu}$ is a \emph{certificate}: a lower bound on the optimal value.
\end{frame}

\board[2]{Weak duality, and the dual of the SVM}{(1) For feasible $\bx$ and $\boldsymbol{\mu} \geq \mathbf{0}$: $d(\boldsymbol{\mu}) \leq \mathcal{L}(\bx, \boldsymbol{\mu}) \leq f(\bx)$. \\ (2) SVM: minimize $\mathcal{L}(\bw, b, \boldsymbol{\mu})$ over $b$ (affine: $\sum_i\mu_iy_i = 0$) and over $\bw$ ($\bw = \sum_i\mu_iy_i\bx_i$), and substitute.}

\checkslide{the dual function and the dual of the SVM}{
    \small
    \textbf{1D example.} $\mathcal{L}(x, \mu) = (x - 2)^2 + \mu(x - 1)$; $\partial_x\mathcal{L} = 2(x - 2) + \mu = 0$ at $x = 2 - \frac{\mu}{2}$:
    \[d(\mu) = \tfrac{\mu^2}{4} + \mu\big(1 - \tfrac{\mu}{2}\big) = \mu - \tfrac{\mu^2}{4}, \qquad d'(\mu) = 1 - \tfrac{\mu}{2} = 0 \text{ at } \mu = 2, \qquad d(2) = 1 = (1 - 2)^2 = p^*.\]
    E.g.\ $d(1) = 0.75 \leq 1$: every $\mu \geq 0$ gives a lower bound.
    \\[0.3em]
    \textbf{SVM.} $\mathcal{L} = \frac{1}{2}\norm{\bw}^2 - \bw^\top\sum_i\mu_iy_i\bx_i - b\sum_i\mu_iy_i + \sum_i\mu_i$: bounded in $b$ only if $\sum_i\mu_iy_i = 0$; minimal in $\bw$ at $\bw = \sum_i\mu_iy_i\bx_i$, where $\frac{1}{2}\norm{\bw}^2 - \bw^\top\bw = -\frac{1}{2}\norm{\bw}^2$:
    \[\max_{\boldsymbol{\mu}}\ \sum_{i=1}^N\mu_i - \frac{1}{2}\sum_{i,j=1}^N\mu_i\mu_jy_iy_j\,\bx_i^\top\bx_j \qquad \text{s.t.} \quad \boldsymbol{\mu} \geq \mathbf{0}, \quad \sum_i\mu_iy_i = 0.\]
    \begin{itemize}
        \item One variable per data point instead of per feature.
        \item The data appear only through the inner products $\bx_i^\top\bx_j$. Replacing them by a kernel $k(\bx_i, \bx_j)$ gives the kernel SVM: a linear classifier in a huge feature space, at the cost of an $N \times N$ Gram matrix.
        \item $\bw = \sum_i\mu_iy_i\bx_i$: only the support vectors contribute.
    \end{itemize}
}

\begin{frame}{Convex conjugates}
    \begin{definition}[][def:conjugate]
        \small $f^*(\by) = \sup_{\bx}\big(\by^\top\bx - f(\bx)\big)$: the largest gap between the line with slope $\by$ and $f$.
    \end{definition}
    \begin{proposition}[][prop:conjugate]
        \small $f^*$ is convex; \ Fenchel--Young: $f(\bx) + f^*(\by) \geq \bx^\top\by$; \ for convex differentiable $f$: $f^*(\nabla f(\bx)) = \bx^\top\nabla f(\bx) - f(\bx)$ and $f^{**} = f$.
    \end{proposition}
    \begin{example}[][ex:conjugates]
        \small
        $\frac{1}{2}\bx^\top Q\bx \leftrightarrow \frac{1}{2}\by^\top Q^{-1}\by$ \ ($\frac{1}{2}\norm{\bx}^2$ is self-conjugate); \ $e^x \leftrightarrow y\log y - y$ ($y > 0$); \ $\norm{\bx}_1 \leftrightarrow 0$ on the $\ell_\infty$ unit ball, $\infty$ outside; \ log-sum-exp $\leftrightarrow \sum_ip_i\log p_i$, the negative entropy (Lecture~10).
    \end{example}
    \small Duals of linear and quadratic programs are again LPs and QPs; a dual feasible point certifies optimality.
\end{frame}

\begin{frame}{Worked examples: computing conjugates}
    \small
    $f^*(y) = \sup_x(yx - f(x))$: maximize $g(x) = yx - f(x)$; for convex differentiable $f$, $g'(x) = 0 \iff f'(x) = y$.
    \\[0.4em]
    \textbf{$f(x) = \frac{1}{2}x^2$:} $y - x = 0$ at $x = y$: \ $f^*(y) = y^2 - \frac{1}{2}y^2 = \frac{1}{2}y^2$. Self-conjugate. (For $\frac{1}{2}\bx^\top Q\bx$: $\bx = Q^{-1}\by$, $f^* = \frac{1}{2}\by^\top Q^{-1}\by$.)
    \\[0.4em]
    \textbf{$f(x) = e^x$:} $y - e^x = 0$ at $x = \log y$ (for $y > 0$): \ $f^*(y) = y\log y - y$. \ $y = 0$: $\sup(-e^x) = 0$. \ $y < 0$: $yx - e^x \to \infty$ as $x \to -\infty$, so $f^*(y) = \infty$.
    \\[0.4em]
    \textbf{$f(x) = |x|$:} $yx - |x| = x(y - \sign x)$. If $|y| \leq 1$, it is $\leq 0$, with $0$ at $x = 0$: $f^*(y) = 0$. If $|y| > 1$, take $x = t\sign y$: $t(|y| - 1) \to \infty$. So $f^*$ is $0$ on $[-1, 1]$ and $\infty$ outside: the $\ell_1$ norm's dual is the $\ell_\infty$ ball.
    \\[0.4em]
    \textbf{Fenchel--Young check} for $e^x$ at $x = 0$, $y = 1$: $f(0) + f^*(1) = 1 + (0 - 1) = 0 = xy$: equality, since $y = f'(0)$.
    \\[0.4em]
    \textbf{An LP certificate.} Minimize $-x_1 - x_2$ s.t.\ $0 \leq x_1, x_2 \leq 1$, $x_1 + x_2 \leq 1.5$. The point $(1, 0.5)$ gives $-1.5$; adding the constraint $x_1 + x_2 \leq 1.5$ with weight $1$ proves $-x_1 - x_2 \geq -1.5$ for every feasible $\bx$: optimal.
\end{frame}

% ------------------------------------------------------------------ Wrap-up
\begin{frame}{Key takeaways}
    \small
    \begin{itemize}
        \item Local min $\Rightarrow \nabla f = \mathbf{0}$; the Hessian's eigenvalues classify critical points.
        \item GD with $\eta \leq \frac{1}{L}$ decreases $f$ and finds small gradients. On quadratics: stable iff $\eta < \frac{2}{\lambda_{\max}}$, rate $\frac{\kappa - 1}{\kappa + 1}$; momentum $\approx\frac{\sqrt\kappa - 1}{\sqrt\kappa + 1}$; Newton: quadratic convergence, expensive.
        \item Convex: every local minimum is global; $\nabla^2 f \succeq 0$. Convex $L$-smooth: $O(1/k)$; $\mu$-strongly convex: $(1 - \mu/L)^k$.
        \item Nonsmooth: $\mathbf{0} \in \partial F(\bx^*)$; proximal gradient; lasso $=$ soft thresholding $\Rightarrow$ sparse weights.
        \item Lagrange: $\nabla f + \sum\lambda_i\nabla g_i = \mathbf{0}$; KKT adds $\mu_i \geq 0$ and $\mu_ig_i = 0$; sufficient for convex problems.
        \item Duality: $d^* \leq p^*$, equality at KKT points of convex problems; the SVM dual needs only inner products (kernels).
    \end{itemize}
\end{frame}

\begin{frame}{Homework and next time}
    \small
    \begin{itemize}
        \item \textbf{Homework:} Lecture 3, Exercises 15, 17, 20, 23, 25 and 27; Programming Exercises P7 and P15.
        \item \textbf{Practice:} Practice 3, sections 9--12.
    \end{itemize}
    \vskip0.6em
    \begin{block}{Next: Probability I--II}
        Probability spaces and Bayes' theorem; random variables and their distributions; expectation, variance and inequalities; joint distributions, covariance, conditional expectation; random vectors and the multivariate Gaussian.
    \end{block}
\end{frame}

\end{document}
